(1) 弦上波速 \(c_i=\sqrt{T/\mu_i}\)，对右行简谐横波的平均功率可写 \(\overline P=\tfrac12T\omega k_i|A|^2=\tfrac12\omega^2Z_i|A|^2\)，其中机械特性阻抗 \(Z_i=T/c_i=\sqrt{T\mu_i}\)。该阻抗不是单纯质量密度，而是波幅与平均能流之间的比例因子。
(2) 设 \(y_1=A(e^{ik_1x}+r e^{-ik_1x})e^{-i\omega t}\)、\(y_2=At e^{ik_2x-i\omega t}\)。接头位移连续给 \(1+r=t\)；同张力且无附加集中质量时斜率连续给 \(k_1(1-r)=k_2t\)。所以 \(r=(k_1-k_2)/(k_1+k_2)=(Z_1-Z_2)/(Z_1+Z_2)\)，\(t=2Z_1/(Z_1+Z_2)\)。
(3) 相对入射功率，反射率 \(\mathcal R=|r|^2=(Z_1-Z_2)^2/(Z_1+Z_2)^2\)；透射功率还应乘阻抗比，\(\mathcal T=(Z_2/Z_1)|t|^2=4Z_1Z_2/(Z_1+Z_2)^2\)。两者相加为一，证明在理想无耗散接头不产生或损失机械能。
(4) 若 \(\mu_2\gg\mu_1\)，\(Z_2\gg Z_1\)，则 \(r\to-1\)，位移反射相位翻转、几乎全反射；若 \(\mu_2\ll\mu_1\)，\(r\to+1\)，同相反射且能量也几乎全返；若 \(Z_1=Z_2\)，则 \(r=0,t=1\)，波无反射传播。后两种极限不可误解为透射幅度和透射能量一定一样。
